% =====================================================================
\chapter{Glossaire}
\label{ann:glossaire}
% =====================================================================

\section{Termes métier}

\begin{description}[leftmargin=0em,style=nextline]
  \item[Souscription (\emph{underwriting}).] Processus par lequel un assureur
  accepte, met à l'étude ou refuse de couvrir un risque, et détermine les
  conditions de la couverture.
  \item[Souscripteur.] Professionnel qui traite les dossiers de souscription au
  quotidien.
  \item[Actuaire.] Professionnel qui construit les modèles de tarification et
  de provisionnement, en s'appuyant sur des données plutôt que sur une
  interface.
  \item[Solvabilité II.] Cadre réglementaire européen de solvabilité des
  assureurs. Il impose notamment la traçabilité et la justification des
  décisions.
  \item[Régime CatNat.] Régime français d'indemnisation des catastrophes
  naturelles, fondé sur la reconnaissance de l'état de catastrophe naturelle
  par arrêté interministériel.
  \item[Non-stationnarité.] Propriété d'un risque dont la distribution change
  dans le temps. Le risque climatique est non stationnaire : le passé ne prédit
  pas directement le futur.
  \item[PPR.] Plan de Prévention des Risques. Document réglementaire qui
  délimite des zones exposées à un aléa et fixe des prescriptions.
  \item[TRI.] Territoire à Risque Important d'inondation, issu de la Directive
  Inondation. Classe l'exposition en quatre niveaux de période de retour.
  \item[Période de retour.] Durée moyenne entre deux occurrences d'un événement
  d'intensité donnée. Une crue centennale a une probabilité d'environ 1\,\% par
  an.
  \item[RGA.] Retrait-gonflement des argiles. Phénomène de déformation des sols
  argileux selon leur humidité, cause majeure de fissuration des bâtiments.
  \item[CatNat historique.] Liste des arrêtés de catastrophe naturelle pris
  pour une commune.
\end{description}

\section{Termes techniques}

\begin{description}[leftmargin=0em,style=nextline]
  \item[WFS.] Web Feature Service. Standard de l'Open Geospatial Consortium
  pour l'échange de données géographiques vectorielles.
  \item[GML.] Geography Markup Language. Format XML d'encodage de géométries,
  utilisé par les services WFS. Le projet traite GML 3.2.
  \item[Point-in-polygon.] Test déterminant si un point est à l'intérieur d'un
  polygone. Implémenté par lancer de rayon.
  \item[Lancer de rayon.] Algorithme consistant à compter les intersections
  d'un rayon partant du point avec les arêtes du polygone. Un nombre impair
  signifie que le point est à l'intérieur.
  \item[Boîte englobante (BBOX).] Rectangle délimitant la zone d'une requête
  spatiale.
  \item[Proximité géométrique.] Test applicable aux géométries non polygonales
  (points, lignes) : le point est-il à moins d'un rayon donné ?
  \item[Ordre des axes.] Convention de notation des coordonnées
  (latitude/longitude ou longitude/latitude). Une erreur d'ordre fait échouer
  silencieusement tout test géométrique.
  \item[Verbatim.] Passage d'une donnée sans transformation : ni renommage, ni
  filtrage, ni recalcul.
  \item[Provenance.] Métadonnée décrivant l'origine d'un fait : fournisseur,
  URL, chaîne d'attribution, horodatage de collecte.
  \item[Résolution.] Niveau auquel un fait a été établi : par bâtiment, par
  commune, ou estimé par commune.
  \item[Mode dégradé.] Comportement du service lorsqu'une source échoue : la
  réponse reste valide mais partielle, avec les faits concernés marqués comme
  inconnus.
  \item[Fixture.] Donnée de test figée, capturée une fois depuis une source
  réelle, permettant des tests hors-ligne reproductibles.
  \item[Serverless.] Modèle d'exécution où le code s'exécute dans une fonction
  éphémère, sans serveur permanent. Le système de fichiers n'y persiste pas.
  \item[Étiquetage.] Mention visible signalant qu'un rendu est illustratif et
  non contractuel.
  \item[Porte de recherche (\emph{grep gate}).] Vérification automatique qui
  échoue si un motif interdit (par exemple un champ de score) apparaît dans le
  code ou un schéma émis.
\end{description}

\section{Sigles}

\begin{table}[H]
\centering
\small
\caption{Sigles et acronymes utilisés dans le rapport.}
\label{tab:sigles}
\begin{tabularx}{\textwidth}{@{}l X@{}}
\toprule
\textbf{Sigle} & \textbf{Signification} \\
\midrule
API & Application Programming Interface \\
BAN & Base Adresse Nationale \\
BDNB & Base de Données Nationale des Bâtiments \\
BBOX & \emph{Bounding box} — boîte englobante \\
CDS & Climate Data Store (Copernicus) \\
CORS & \emph{Cross-Origin Resource Sharing} \\
CSTB & Centre Scientifique et Technique du Bâtiment \\
DPE & Diagnostic de Performance Énergétique \\
GML & Geography Markup Language \\
IGN & Institut Géographique National \\
INSEE & Institut National de la Statistique et des Études Économiques \\
LLM & \emph{Large Language Model} — grand modèle de langage \\
OGC & Open Geospatial Consortium \\
PPR & Plan de Prévention des Risques \\
PPRI & Plan de Prévention du Risque Inondation \\
RCP & \emph{Representative Concentration Pathway} — trajectoire de
  concentration représentative (scénario climatique) \\
REST & \emph{Representational State Transfer} \\
RGA & Retrait-gonflement des argiles \\
SSP & Sites et sols pollués \\
SPA & \emph{Single Page Application} \\
TMD & Transport de matières dangereuses \\
TRI & Territoire à Risque Important d'inondation \\
TTL & \emph{Time To Live} — durée de vie (d'un cache) \\
VFX & \emph{Visual effects} — effets visuels (mode de rendu cinématique) \\
WFS & Web Feature Service \\
WMS & Web Map Service \\
\bottomrule
\end{tabularx}
\end{table}
